\documentclass[a4paper,10pt]{report}\input{../0.Préambule/Préambule_cours_prof}\begin{document}

\pagestyle{empty}
\newpage
\noindent NOM, Prénom et classe : \dotfill

\section*{Enchainement d'opérations \hspace{3cm} Sujet A}
\addcontentsline{toc}{section}{Enchainement d'opérations \hspace{3cm} Sujet A}

\noindent \textit{L'usage de la calculatrice est \textbf{interdite}.} \hfill \textit{Durée : 30 minutes}

\paragraph{Exercice 1} : \quad Qu'est ce qu'une différence ? \hfill \textbf{2 points}

\paragraph{Exercice 2} : \quad Entoure le signe opératoire de l'opération prioritaire. (Il peut y en avoir plusieurs) \hfill \textbf{4 points}

\begin{enumerate}[font=\color{exer1}]
\begin{tabularx}{\linewidth}{*{4}{>{\arraybackslash}X}}
\item $S = 9 \div 3 + (15 - 6 \div 3)$ &
\item $T = 34 - ( 104 \div 52 \times 5)$ &
\item $O = (84 - 1) \div (5 + 0,4)$ &
\item $P = 2 \times \left( (1+2) \times 4 - 2 \right)$ \\
\end{tabularx}
\end{enumerate}

\paragraph{Exercice 3} : \quad Effectue les calculs en soulignant le calcul en cours.\hfill \textbf{9 points}

\begin{enumerate}[font=\color{exer1}]
\begin{tabularx}{\linewidth}{*{3}{>{\arraybackslash}X}}
\item $M = 18 - \left( 4 \times (5 - 4) + 2 \right)$ &
\item $E = 24 \div \left(8 - (4 + 1) \right)$ &
\item $T = \left( 2 + 2 \times (5 + 4) \right) \div 4$ \\
\item $S = \dfrac{17 - 5}{3} + 2$ &
\item $V = \dfrac{45,5}{2} \times 3 - 1$ &
\item $I = \dfrac{27}{2 \times 3} - 1$ \\
\end{tabularx}
\end{enumerate}

\paragraph{Exercice 4} : \hfill \textbf{3 points}

\noindent \textit{La rédaction sera évaluée au même titre que les résultats, on prendra donc soin d'écrire chacun de nos calculs.}

\noindent Maya achète $5$ pots de confitures à $1,80$\euro pièce et $12$ baguettes de pain à $0,70$\euro pièce. 

\noindent Quel est le prix total qu'elle doit payer ?

\paragraph{Exercice 5} : \quad Calcul le produit de la somme de $5$ et de $2$ par $3$ \hfill \textbf{2 points}

\vspace{3cm}
\noindent NOM, Prénom et classe : \dotfill

\section*{Enchainement d'opérations \hspace{3cm} Sujet B}
\addcontentsline{toc}{section}{Enchainement d'opérations \hspace{3cm} Sujet B}

\noindent \textit{L'usage de la calculatrice est \textbf{interdite}.} \hfill \textit{Durée : 30 minutes}

\paragraph{Exercice 1} : \quad Qu'est ce qu'un quotient ? \hfill \textbf{2 points}

\paragraph{Exercice 2} : \quad Entoure le signe opératoire de l'opération prioritaire. (Il peut y en avoir plusieurs) \hfill \textbf{4 points}

\begin{enumerate}[font=\color{exer1}]
\begin{tabularx}{\linewidth}{*{4}{>{\arraybackslash}X}}
\item $S = 9 \div 3 + (15 - 6 \div 2)$ &
\item $P = 3 \times \left( (1+2) \times 4 - 2 \right)$&
\item $O = (84 - 1) \div (5 + 0,4)$ &
\item $T = 34 - ( 104 \div 52 \times 6)$  \\
\end{tabularx}
\end{enumerate}

\paragraph{Exercice 3} : \quad Effectue les calculs en soulignant le calcul en cours.\hfill \textbf{9 points}

\begin{enumerate}[font=\color{exer1}]
\begin{tabularx}{\linewidth}{*{3}{>{\arraybackslash}X}}
\item $M = 18 - \left( 4 \times (5 - 2) + 2 \right)$ &
\item $T = 24 \div \left(8 - (3 + 1) \right)$ &
\item $S = \left( 2 + 2 \times (5 + 2) \right) \div 4$ \\
\item $A = \dfrac{81}{9} \times 5 - 1$ &
\item $E = 7 \times \dfrac{15 \times 4}{3 - 2} + 2 \times 8$ &
\item $Z = \dfrac{13 \times (4 + 7) - 5}{13 - (2 \times 4 + 3}$ \\
\end{tabularx}
\end{enumerate}

\paragraph{Exercice 4} : \hfill \textbf{3 points}

\noindent \textit{La rédaction sera évaluée au même titre que les résultats, on prendra donc soin d'écrire chacun de nos calculs.}

\noindent Érine achète $5$ pots de confitures à $1,80$\euro pièce et $11$ baguettes de pain à $0,70$\euro pièce. 

\noindent Quel est le prix total qu'elle doit payer ?

\paragraph{Exercice 5} : \quad Calcul la somme de $5$ et du produit de $2$ par $3$ \hfill \textbf{2 points}
\end{document}
